\documentclass[11pt,a4paper]{article}
\usepackage{fontspec}
\setmainfont{Latin Modern Roman}
\usepackage[margin=23mm]{geometry}
\usepackage{amsmath,amssymb,amsthm,mathtools,mathrsfs}
\usepackage{longtable,booktabs}
\usepackage[unicode,hidelinks]{hyperref}
\providecommand{\C}{\mathbb C}
\providecommand{\R}{\mathbb R}
\newtheorem{theorem}{Theorem}[section]
\newtheorem{lemma}[theorem]{Lemma}
\newtheorem{proposition}[theorem]{Proposition}
\newtheorem{corollary}[theorem]{Corollary}
\newtheorem{definition}[theorem]{Definition}
\newtheorem{remark}[theorem]{Remark}
\allowdisplaybreaks
\emergencystretch=3em
\makeatletter
\renewcommand\@pnumwidth{2.5em}
\renewcommand\@tocrmarg{3.5em}
\makeatother
\title{The Actual Period Quadric and Residue Commutator}
\author{Split-Zero research programme}
\date{14 September 2026}
\begin{document}
\maketitle
\paragraph{Portable source edition.}
Every TeX input used by this entry is expanded below. Historical file
and directory names in the preserved text are source locators; the
exact origin paths and hashes are recorded in \texttt{11\_SOURCE\_MAP.json}.
The other numbered proof files in this flat handoff supply the
accompanying complete calculations.\par

\section{The actual period quadric and its cyclic continuation}
\label{sec:PQO}

This calculation uses the original simple quartet, $m=1$, and the
literal period matrix, amplitude unit and cyclic matrix of BRU and
OPG. Its formulas evaluate the transported quadric from those data.
All source orders elsewhere in the programme remain unchanged.

\subsection{Original roots, entire unit and centred phase}
Order the roots and keep their actual values as
\[
\begin{gathered}
 \alpha_1=\tfrac12+\delta+i\gamma,\quad
 \alpha_2=\tfrac12+\delta-i\gamma,\quad
 \alpha_3=\tfrac12-\delta+i\gamma,\quad
 \alpha_4=\tfrac12-\delta-i\gamma,\\
 0<\delta<\tfrac12,\quad\gamma>2,\quad
 h(s)=\prod_{j=1}^4(s-\alpha_j),\quad g=2\xi,\quad v_h=g/h.
\end{gathered}\tag{PQO1}
\]
The four zero orders of $g$ are complete and equal to one.
Thus $v_h$ is entire and has nonzero value at each displayed root.
Put $w=s-1/2$, and define the following original constants:
\[
\begin{gathered}
 z_+=\delta+i\gamma,\quad z_-=\delta-i\gamma,\quad
 \lambda=z_+^2,\quad\mu=z_-^2,\\
 A=2(\gamma^2-\delta^2),\qquad B=(\delta^2+\gamma^2)^2,\qquad
 \mathfrak a=\Phi_h(1/2)=\frac1{160}+\frac A{24}+\frac B2,\\
 h(1/2+w)=w^4+Aw^2+B,\qquad
 \Phi_h(1/2+w)=\mathfrak a+\frac{w^5}{5}+\frac{Aw^3}{3}+Bw.
\end{gathered}\tag{PQO2}
\]
The product of the two quadratics
$((w-\delta)^2+\gamma^2)((w+\delta)^2+\gamma^2)$ proves the
quartic identity. Integrating it and using $\Phi_h(0)=0$
proves the phase and constant in PQO2. In particular the translation
has not removed the original constant $\mathfrak a$.

The theta integral IPM.2 has real coefficients, real theta kernel,
and an expression unchanged by $s\mapsto1-s$. It therefore gives
$g(\bar s)=\overline{g(s)}$ and $g(1-s)=g(s)$ directly; the same
two identities hold for $h$. Taking their entire quotient gives
\[
 v_h(\alpha_1)=v_h(\alpha_4)=a,\qquad
 v_h(\alpha_2)=v_h(\alpha_3)=\bar a,\qquad a\ne0.
 \tag{PQO3}
\]
These are the full actual unit values. Their magnitude and phase
are not assigned. Set $t_+=a^{-2}$ and $t_-=\bar a^{-2}$.
For $j=0,1,2,3$ retain
\[
 \sigma_j=t_+\lambda^j-t_-\mu^j,\qquad
 \sigma_2=-A\sigma_1-B\sigma_0,\qquad
 \sigma_3=(A^2-B)\sigma_1+AB\sigma_0.
 \tag{PQO4}
\]
The recurrence follows because $\lambda+\mu=-A$ and
$\lambda\mu=B$. Each $\sigma_j$ is purely imaginary.
Furthermore $\sigma_0$ and $\sigma_1$ cannot both vanish:
the first would give $t_+=t_-\ne0$, and the second would then
give $t_+(\lambda-\mu)=0$, whereas
$\lambda-\mu=4i\delta\gamma\ne0$.

\subsection{The original periods as an absolutely convergent series}
Fix $u\ne0$ on its original logarithm branch. Let
$\ell_j(r)=r\exp(i(\arg u+\pi+2\pi j)/5)$,
$\Gamma_j=-\ell_0+\ell_j$, $1\leq j\leq4$.
Let $\Pi$ integrate the degree-less-than-four representative
against $e^{\Phi_h(s)/u}\,ds$ on these original oriented contours.
Write $F_{ja}=\zeta^{ja}-1$, $\zeta=e^{2\pi i/5}$,
$1\leq j,a\leq4$. Direct finite geometric sums prove
\[
 (F^{-1})_{aj}=\frac15\zeta^{-ja},\qquad
 CF=F D,\qquad D=\operatorname{diag}(\zeta,\zeta^2,\zeta^3,\zeta^4),
 \tag{PQO5}
\]
where $(Cx)_j=x_{j+1}-x_1$ for $j<4$ and $(Cx)_4=-x_1$.
For example multiplying the asserted inverse by $F$ gives
$\frac15\sum_{j=1}^4\zeta^{-ja}(\zeta^{jb}-1)=\delta_{ab}$;
substitution into each row gives $CF=FD$.

Let $T$ be the exact increasing-coefficient translation matrix
$T_{rb}=\binom br2^{-(b-r)}$ for $0\leq r\leq b\leq3$, and
zero otherwise. Thus it sends coefficients of $p(s)$ to those
of $p(1/2+w)$, with determinant one. Define the centred
Fourier-period matrix
\[
 P_{ar}(u)=\frac15\sum_{j=1}^4\zeta^{-ja}
 \int_{\Gamma_j}e^{(w^5/5+Aw^3/3+Bw)/u}w^r\,dw,
 \quad 1\leq a\leq4,\quad0\leq r\leq3.
 \tag{PQO6}
\]
The contours in this formula start at $w=0$. They give precisely
the translated original period by
\[
 F^{-1}\Pi=e^{\mathfrak a/u}P(u)T.
 \tag{PQO7}
\]
To justify this exact contour translation, the substitution
$w=s-1/2$ first gives rays starting at $-1/2$. Join each shifted
ray at a large radius to the corresponding unshifted ray. On the
joining segment their arguments differ by $O(1/r)$, and the
leading real phase is $-r^5/(5|u|)+O(r^4)$; all lower phase
terms and polynomial factors are dominated by this negative
leading term. Those outer joining integrals tend to zero.
The two common-origin joining integrals cancel in
$-\ell_0+\ell_j$ with their opposite orientations. Cauchy's
theorem for the entire integrand proves equality in the limit.
Substitution of the polynomial gives exactly the matrix $T$.

For every integer $L\geq1$ define, on the same branch,
\[
 G_L(u)=u^{L/5}5^{L/5-1}e^{\pi iL/5}\Gamma(L/5).
 \tag{PQO8}
\]
Integration of the leading phase on the two rays gives
$\int_{\Gamma_j}e^{w^5/(5u)}w^{L-1}dw
 =G_L(u)(\zeta^{jL}-1)$ by the real substitution
$r^5/(5|u|)$ on each ray, including its phase in $dw$.
Expanding the remaining two exponentials now gives the exact
entries
\[
\boxed{
 P_{ar}(u)=
 \sum_{\substack{p,q\geq0\\r+3p+q+1\equiv a\pmod5}}
 \frac{A^pB^q}{3^p p!q!u^{p+q}}\,
                 G_{r+3p+q+1}(u).
}\tag{PQO9}
\]
Terms with $r+3p+q+1\equiv0$ have zero two-ray integral
in every column and contribute zero before the Fourier map.
To prove every exchange, the sum of the absolute values of the
unintegrated series on either ray is bounded by a constant times
\[
 R^r\exp\left(-\frac{R^5}{5|u|}
             +\frac{|A|R^3}{3|u|}+\frac{|B|R}{|u|}\right).
\]
Here $R$ is the radial coordinate and $r$ is the fixed column
degree. This integrable majorant
proves absolute termwise integration. It is uniform on compact
coefficient and nonzero-$u$ sets on the chosen branch. The finite
Fourier sum in PQO5 then selects exactly the stated residue class.
Thus PQO9 is an actual period formula, with all coefficients and
all terms retained.

For additional finite computation, let $I_d$ denote the vector
of the four centred ray moments of $w^d$. Integration of the
derivative of $e^{(w^5/5+Aw^3/3+Bw)/u}w^d$, with the same
vanishing outer endpoints and cancelling inner endpoints, proves
\[
 I_{d+4}+A I_{d+2}+B I_d=-ud I_{d-1}\quad(d\geq0),
 \qquad u\partial_B I_d=I_{d+1},\quad
 3u\partial_A I_d=I_{d+3}.
 \tag{PQO10}
\]
The first right side is zero at $d=0$. Coefficient derivatives
are justified by the same compact domination. These are the
original polynomial period recurrences, and they compute every
higher moment from the first four and the actual coefficients.

The determinant is especially explicit in these centred
coordinates:
\[
 \det\Pi=-(2\pi)^2u^2e^{4\mathfrak a/u},\qquad
 \det F=-25\sqrt5,\qquad
 \boxed{\det P(u)=\frac{(2\pi)^2}{25\sqrt5}\,u^2.}
 \tag{PQO11}
\]
For the first equality use the full PDE1--13 determinant
calculation at $n=4$: the phase is $e^{15\pi i}=-1$ and
the critical-value sum is $4\mathfrak a$. This critical sum
also follows directly by pairing $w$ with $-w$ in the odd
centred phase. For $F$, subtract row zero from the other rows
in the five-by-five Fourier Vandermonde and expand its first
column. Its phase is $e^{13\pi i}=-1$ and its magnitude
$5^{5/2}$, by the finite Fourier orthogonality identity.
The last equality follows by taking determinants in PQO7 and
using $\det T=1$. The Gamma grid-product proof PGCC1--6 retains
the full evaluated constant in the first equality. In particular
every inverse below exists on the full stated domain.

\subsection{The actual transformed quadric and every deciding coefficient}
Let $V$ be the centred evaluation matrix with rows
$(1,z_j,z_j^2,z_j^3)$, where
$(z_1,z_2,z_3,z_4)=(z_+,z_-,-z_-,-z_+)$. Its inverse sends
the four primary values to their unique degree-less-than-four
polynomial. Thus the actual matrix from primary coordinates
to the Fourier-period space is
\[
 B_{\rm per}=F^{-1}\Pi UJ_{\rm prim}
       =e^{\mathfrak a/u}P(u)V^{-1}
                         \operatorname{diag}(a,\bar a,\bar a,a).
 \tag{PQO12}
\]
This follows either by applying all maps to the four CRT
idempotents or by evaluating the full unit at each root.
It agrees with OCS and OPG in the given original orientation.

For $Q_0(t)=t_1t_4-t_2t_3$, define
$Q_u(x)=Q_0(B_{\rm per}^{-1}x)$. On the unchanged original
period coordinate $\mathbf y\in\mathbb C^4$ the quadratic is
$Q_u^{\rm per}(\mathbf y)=Q_0((\Pi UJ_{\rm prim})^{-1}\mathbf y)
 =Q_u(F^{-1}\mathbf y)$, because $FB_{\rm per}=\Pi UJ_{\rm prim}$.
Thus the exact isomorphism $\mathbf y=Fx$ carries the following
quadric to the original period space. If
$p(w)=p_0+p_1w+p_2w^2+p_3w^3$, multiplication gives
\[
 p(w)p(-w)=p_0^2+(2p_0p_2-p_1^2)w^2
                     +(p_2^2-2p_1p_3)w^4-p_3^2w^6.
\]
Consequently the exact symmetric matrix in the centred
coefficient frame is
\[
 S=\begin{pmatrix}
 \sigma_0&0&\sigma_1&0\\
 0&-\sigma_1&0&-\sigma_2\\
 \sigma_1&0&\sigma_2&0\\
 0&-\sigma_2&0&-\sigma_3
 \end{pmatrix},\qquad
 \boxed{Q_u(x)=x^{\mathsf T}H(u)x,
 \quad H(u)=e^{-2\mathfrak a/u}P(u)^{-\mathsf T}S P(u)^{-1}.}
 \tag{PQO13}
\]
Indeed applying the preceding product formula at $z_+$ and
$z_-$ with multipliers $t_+$ and $-t_-$ gives exactly $S$.
This retains both actual unit values through every $\sigma_j$.
The two parity blocks have determinants
\[
 \sigma_0\sigma_2-\sigma_1^2
          =-t_+t_-(\lambda-\mu)^2,\qquad
 \sigma_1\sigma_3-\sigma_2^2
          =-t_+t_-B(\lambda-\mu)^2.
\]
Thus $S$ and $H(u)$ are nondegenerate, with the exact determinant
\[
 \det H(u)=e^{-8\mathfrak a/u}
       \frac{B(t_+t_-)^2(\lambda-\mu)^4}{(\det P(u))^2}\ne0.
 \tag{PQO14}
\]

Write $N(u)=\operatorname{adj}P(u)$, with entries $N_{ra}$,
$0\leq r\leq3$, $1\leq a\leq4$, and retain
$\Delta(u)=\det P(u)$ from PQO11. Then every coefficient is
given by the explicit minor formula
\[
\begin{aligned}
 H_{ab}(u)=\frac{e^{-2\mathfrak a/u}}{\Delta(u)^2}
 \{&\sigma_0N_{0a}N_{0b}
 +\sigma_1(N_{0a}N_{2b}+N_{2a}N_{0b}-N_{1a}N_{1b})\\
 &+\sigma_2(N_{2a}N_{2b}-N_{1a}N_{3b}-N_{3a}N_{1b})
                         -\sigma_3N_{3a}N_{3b}\}.
\end{aligned}\tag{PQO15}
\]
This is matrix multiplication in PQO13 using
$P^{-1}=N/\Delta$. Each $N_{ra}$ is the signed original
three-by-three period minor. Its entries are given by the
absolutely convergent series PQO9 or by the original integrals
PQO6. The polynomial coefficient of $x_a^2$ is $H_{aa}$;
that of $x_ax_b$, $a<b$, is $2H_{ab}$. Both factors are
retained in this specified quadratic form.

The exact eight deciding entries are
\[
 \mathcal E(u)=
 (H_{11},H_{22},H_{33},H_{44},H_{12},H_{13},H_{24},H_{34}).
 \tag{PQO16}
\]
The cyclic action on quadrics is the pullback
$Q_u\mapsto Q_u(D^{-1}x)$; its symmetric matrix is
$D^{-\mathsf T}H(u)D^{-1}$. Since $D^5=I$ and $\det D=1$,
the line spanned by a nondegenerate form can be stable only
with character one. In fact if its multiplier is $\eta$, then
$\eta^5=1$ by the group action, while taking determinants
gives $\eta^4=1$ because $\det H(u)\ne0$. Thus $\eta=1$.
Entrywise, invariance is precisely
$(\zeta^{-(a+b)}-1)H_{ab}=0$. The only pairs with
$a+b\equiv0$ are $(1,4)$ and $(2,3)$. Therefore
\[
\boxed{\begin{gathered}
 \text{The quadric is invariant exactly on }
       Z=\{u:\mathcal E(u)=0\};\\
 u\in Z\ \Longrightarrow\
 Q_u(x)=2H_{14}(u)x_1x_4+2H_{23}(u)x_2x_3,
                 \quad H_{14}(u)H_{23}(u)\ne0;\\
 \text{at every }u\notin Z\text{ its orbit consists of five distinct quadrics.}
\end{gathered}}\tag{PQO17}
\]
Every assertion follows from the entry test and the rank-four
determinant. The orbit size divides five; it is one exactly
when the line is stable. Thus no two-member orbit occurs.
Moreover $F^{-1}C^{-1}=D^{-1}F^{-1}$ by PQO5, so this is
exactly the orbit of $Q_u^{\rm per}(C^{-j}\mathbf y)$ in the
original period coordinates, with the same five indexed maps.
PQO15--17 determine the exact exceptional locus by original
period minors and original amplitude values, without assigning
generic values to them.

\subsection{The cyclic action is the actual original continuation in $u$}
Continue the original logarithm of $u$ by $2\pi i$. Its rays
advance from $\ell_j$ to $\ell_{j+1}$, and their oriented
differences advance by the exact original matrix $C$.
For the last row $\ell_5=\ell_0$, giving $-\Gamma_1$.
Equivalently every summand of PQO9 acquires the phase
\[
 e^{2\pi i((r+3p+q+1)/5-p-q)}
       =\zeta^{r+3p+q+1}=\zeta^a.
\]
The original constant $e^{\mathfrak a/u}$ is single-valued
under this continuation. Hence
\[
 P(\log u+2\pi i)=DP(\log u),\quad
 B_{\rm per}(\log u+2\pi i)=DB_{\rm per}(\log u),\quad
 H(\log u+2\pi i)=D^{-\mathsf T}H(\log u)D^{-1}.
 \tag{PQO18}
\]
The contour argument and the convergent series prove the same
map. Thus PQO17 describes the actual period continuation of
this original quadric. It introduces no freely chosen cyclic
action on its coefficients.

\subsection{The large-parameter geometry and a positive-ray obstruction}
Let $\varepsilon=u^{-1/5}$ on the same branch, and set
$D_\varepsilon=\operatorname{diag}(\varepsilon^{-1},
\varepsilon^{-2},\varepsilon^{-3},\varepsilon^{-4})$.
The exact series PQO9 gives
\[
 P(u)=\mathscr P(A\varepsilon^2,B\varepsilon^4)D_\varepsilon,
 \qquad
 \mathscr P(0,0)=\operatorname{diag}(g_1,g_2,g_3,g_4),
 \quad g_a=5^{a/5-1}e^{\pi ia/5}\Gamma(a/5).
 \tag{PQO19}
\]
The entries of $\mathscr P$ are entire in the two displayed
arguments, by the same absolute integrable series majorant.
The same determinant proof PQO11 applies to every complex pair
of lower odd-phase coefficients, since the centred phase remains
odd and its critical values pair to zero with full multiplicities.
Thus their determinant is the nonzero constant
$(2\pi)^2/(25\sqrt5)$, by PQO11 and $\det D_\varepsilon=u^2$.
Their inverse is therefore holomorphic everywhere, given by its
adjugate divided by that precise constant. Near $\varepsilon=0$
it is $\operatorname{diag}(g_a^{-1})+O(\varepsilon^2)$.
This bound follows from its convergent power series after
substituting $(A\varepsilon^2,B\varepsilon^4)$; it holds on
a full small complex disc, with the fixed original $A,B$.

If $\sigma_0\ne0$, PQO13 and PQO19 give
\[
 e^{2\mathfrak a/u}H_{11}(u)
              =\frac{\sigma_0}{g_1^2}\varepsilon^2
                                     +O(\varepsilon^4).
\]
If $\sigma_0=0$, then $\sigma_1\ne0$ by PQO4, and they give
\[
 e^{2\mathfrak a/u}H_{22}(u)
              =-\frac{\sigma_1}{g_2^2}\varepsilon^4
                                     +O(\varepsilon^6).
 \tag{PQO20}
\]
For verification of the orders, insert
$P^{-1}=D_\varepsilon^{-1}\mathscr P^{-1}$ into PQO13.
The middle matrix has entries $S_{rs}\varepsilon^{r+s+2}$.
Its first possible degree is two, from $S_{00}$, and its
next possible degree is four, from $S_{02},S_{20},S_{11}$.
The constant term of $\mathscr P^{-1}$ is diagonal, so the
indicated leading entry is exactly the displayed nonzero
coefficient. Every further contribution has the stated higher
degree. This also proves the expansion when some other
$\sigma_j$ vanish.

Consequently for every fixed original quartet and its full
actual amplitude there exists a finite radius outside which
the actual orbit has five distinct quadrics. This conclusion
holds on every original logarithm branch: the estimate is on
the full complex $\varepsilon$ disc, and its leading coefficient
is nonzero. The set $Z$ of PQO17 is a discrete subset of the
logarithmic $u$ surface. Indeed the decisive entry just proved
is holomorphic on that connected surface and is not identically
zero; its zeros are isolated by its local power series, and
$Z$ is contained in that zero set. The explicit finite-radius
bound from the companion period-connection calculation
strengthens this argument while retaining the same actual
coefficients.

Moreover PQO18 multiplies each deciding entry by its nonzero
root-of-unity factor, so the condition $\mathcal E(u)=0$ is
unchanged by every full logarithm turn. Therefore $Z$ is the
inverse image under $\log u\mapsto u$ of a well-defined set
$Z_{\rm abs}\subset\mathbb C^*$. On a sufficiently small disc
about any nonzero $u$, choose one logarithm branch. The preceding
isolated-zero proof shows that $Z_{\rm abs}$ is discrete and
closed on that disc. The explicit radius in the companion
calculation gives $|u|<R_*$ throughout $Z_{\rm abs}$.
Consequently its only possible accumulation point in the finite
$u$-plane is zero. In particular each compact annulus about
zero contains only finitely many such points: an infinite
subset of that compact set would have a nonzero accumulation
point, contradicting the proved local isolated-zero property.

For positive real $u$ retain the actual lift
$\log u=\ln u+2\pi i\ell_{\rm br}$, $\ell_{\rm br}\in\mathbb Z$.
The exact series has an additional real structure. Every
surviving index in PQO9 is $a+5L$, and its phase at the lift
$\ell_{\rm br}=0$ is $e^{\pi ia/5}(-1)^L$. Its remaining
coefficients are real. The actual lift multiplies row $a$ by
$\zeta^{\ell_{\rm br}a}$, by PQO18. Absolute convergence
therefore proves
\[
 P(u)=E_{\ell_{\rm br}} R(u),\qquad
 E_{\ell_{\rm br}}=D^{\ell_{\rm br}}
 \operatorname{diag}(e^{\pi i/5},e^{2\pi i/5},
                       e^{3\pi i/5},e^{4\pi i/5}),
 \qquad R(u)\text{ real and invertible}.
 \tag{PQO21}
\]
Let $M$ be multiplication by $w$ in the increasing centred
polynomial frame, so
\[
 M=\begin{pmatrix}0&0&0&-B\\1&0&0&0\\0&1&0&-A\\0&0&1&0\end{pmatrix},
 \quad c_3=\frac{t_+-t_-}{4i\delta\gamma},\quad
 c_1=\frac{t_++t_-}{2}-(\delta^2-\gamma^2)c_3.
 \tag{PQO22}
\]
Both $c_1,c_3$ are real by PQO3, and the polynomial
$c_1+c_3 w^2$ takes values $t_+,t_-,t_-,t_+$ at the four
original roots. Hence $c_1M+c_3M^3$ is precisely multiplication
by $w(j_hv_h)^{-2}$, with its entire unit retained.

The companion period-connection proof establishes the exact
identity between cyclic invariance of $Q_u$ and
\[
 [P(u)(c_1M+c_3M^3)P(u)^{-1},D]=0.
 \tag{PQO23}
\]
It proves this by the flat residue/reflection form on the
original polynomial periods and retains all its constants.
Using that proved identity, PQO21 gives a further precise
obstruction on the positive ray. Since $D$ has distinct
diagonal entries, a matrix commuting with it is diagonal.
The matrix in PQO23 is $E_{\ell_{\rm br}}$ times a real matrix
times $E_{\ell_{\rm br}}^{-1}$;
if diagonal, its diagonal entries must therefore be real.
Its eigenvalues are exactly
\[
 \frac{\delta+i\gamma}{a^2},\quad
 \frac{\delta-i\gamma}{\bar a^2},\quad
 -\frac{\delta-i\gamma}{\bar a^2},\quad
 -\frac{\delta+i\gamma}{a^2},
 \tag{PQO24}
\]
as follows by evaluation at the four distinct roots of $h$.
Thus every positive-real-$u$ point in $Z$ satisfies the explicit
original-unit equation
\[
 \operatorname{Im}\left(\frac{\delta+i\gamma}{a^2}\right)=0.
 \tag{PQO25}
\]
If that imaginary part is nonzero for the original unit, the
positive ray contains no point of $Z$. When it is zero, the
four displayed eigenvalues are the two nonzero real numbers
$r,-r$, each twice, and the complete period-minor tests
PQO15--17 still determine $Z$. The equation is not assumed
to hold or fail. Together the exact minor tests, the actual
period continuation, the proved large-parameter result and
this real-structure obstruction give control of the original
quadric family with every amplitude value and phase present.


% Independent mathematical provider. Original periods, amplitude and branch retained.
\section{The original period quadric as a residue commutator}
\label{sec:actual-period-residue-commutator}

Retain $0<\delta<1/2$, $\gamma>2$, the ordered simple quartet
$\alpha_1=1/2+\delta+i\gamma$, $\alpha_2=1/2+\delta-i\gamma$,
$\alpha_3=1/2-\delta+i\gamma$, $\alpha_4=1/2-\delta-i\gamma$,
and $h(s)=\prod_a(s-\alpha_a)$. Put
\[
 w=s-\tfrac12,\quad r_a=\alpha_a-\tfrac12,\quad
 \beta=2(\gamma^2-\delta^2),\quad \eta=(\delta^2+\gamma^2)^2,
 \quad H(w)=w^4+\beta w^2+\eta.
 \tag{RPC1}
\]
Thus $r_4=-r_1$, $r_3=-r_2$, $r_2=\overline{r_1}$ and
$r_1^2-r_2^2=4i\delta\gamma$. No root or original coordinate is
discarded by this translation. Its coefficient isomorphism is
$J_c:\mathbb C[w]_{<4}\to E_h$, $J_cp=[p(s-1/2)]$; its inverse
is substitution $s=w+1/2$ in the unique original representative.
The primitive with its original constant is
\[
 \Phi_h(s)=A+\phi(w),\quad
 \phi(w)=\frac{w^5}{5}+\frac{\beta w^3}{3}+\eta w,\quad
 A=\Phi_h(1/2)=\frac1{160}+\frac\beta{24}+\frac\eta2.
 \tag{RPC2}
\]
The last equality is forced by $\Phi_h(0)=0$ and follows by evaluating
the displayed odd polynomial at $w=-1/2$.

Let $u\ne0$ carry the specified original branch of its logarithm.
Let $\Pi$ be the original period map on the rays and orientations
of PDE2, and let $F_{ja}=\zeta^{ja}-1$, $\zeta=e^{2\pi i/5}$,
$1\leq j,a\leq4$. Its inverse has entries
$(F^{-1})_{aj}=\zeta^{-ja}/5$: multiplying the two matrices and using
$\sum_{j=1}^4\zeta^{jb}=4$ for $5\mid b$ and $-1$ otherwise proves
this entry by entry. The literal centered coefficient-to-Fourier map is
\[
 P=F^{-1}\Pi J_c,\qquad P_0=e^{-A/u}P.
 \tag{RPC3}
\]
These are invertible by the complete original period determinant
PDE3--13. The contours for $P_0$ may be taken to be the rays based at
$w=0$ with the original five angles. Indeed translating the old rays
gives rays based at $-1/2$. Join corresponding finite endpoints and
close each contour at radius $R$. The integrand is entire. On the
bounded-length outer joining curves its exponent is
$-R^5/(5|u|)+O(R^4)$, uniformly along the joining curves, so the
outer integrals tend to zero. The common finite joining path occurs
with opposite coefficients in $-\ell_0+\ell_j$ and cancels exactly.
Cauchy's integral theorem therefore proves this contour identity,
including its signs and the factor $e^{A/u}$.

\subsection{A flat alternating pairing with all constants}

On $\mathbb C[w]/(H)$ define the bilinear, rather than sesquilinear,
pairing
\[
 \mathcal J(p,q)=\sum_{H(r)=0}\frac{p(r)q(-r)}{H'(r)}.
 \tag{RPC4}
\]
The roots are simple at the original packet. The equivalent polynomial
formula is the coefficient of $w^3$ in the remainder of $p(w)q(-w)$
modulo $H$. To prove it, use Lagrange interpolation: the coefficient
of $w^3$ in the idempotent at $r$ is $1/H'(r)$. This polynomial
formula extends the pairing to every complex value of $\beta,\eta$,
including repeated roots, without a division by a repeated-root value.
In the increasing frame $1,w,w^2,w^3$, multiplication by $w$ and
this pairing are the exact matrices
\[
 M=\begin{pmatrix}0&0&0&-\eta\\1&0&0&0\\0&1&0&-\beta\\0&0&1&0\end{pmatrix},
 \qquad
 J=\begin{pmatrix}0&0&0&-1\\0&0&1&0\\0&-1&0&\beta\\1&0&-\beta&0\end{pmatrix}.
 \tag{RPC5}
\]
For the last matrix, the coefficient functional is zero through degree
two, one in degree three, zero in degree four and six, and $-\beta$
in degree five. Inserting the sign $(-1)^j$ from the second argument
gives every displayed entry. Thus $J^{\mathsf T}=-J$, $\det J=1$,
and direct multiplication gives $M^{\mathsf T}J+JM=0$.

Allow the two original odd-phase coefficients to vary complexly,
solely to calculate the existing period map. The ray domination
$e^{-r^5/(5|u|)+O(r^3+r)}$ on compact parameter sets justifies all
coefficient derivatives and integrations by parts below. Ordinary
division of $w^{j+3}$ by $H$, followed by differential reduction,
gives
\[
 \partial_\eta P_0=\frac1uP_0M,\qquad
 \partial_\beta P_0=\frac1{3u}P_0(M^3-uN),\qquad
 N=\begin{pmatrix}0&0&1&0\\0&0&0&2\\0&0&0&0\\0&0&0&0\end{pmatrix}.
 \tag{RPC6}
\]
For clarity, the quotient derivatives for columns $j=0,1,2,3$ are
$0,0,1,2w$, respectively; hence this is the full correction matrix.
The two outer boundary terms vanish, and the two values at the common
origin cancel. No differential remainder was replaced by an arithmetic
remainder in RPC6.

There is also a direct nonvanishing proof on this entire two-parameter
family. The displayed matrices satisfy $\operatorname{Tr}M=0$,
$\operatorname{Tr}M^3=0$ and $\operatorname{Tr}N=0$. The first and
third follow from their diagonals; the second follows by multiplying
RPC5 three times, or from the vanishing third Newton sum of the even
monic polynomial $H$. Multilinearity of the determinant in its columns
then gives $\partial_\eta\det P_0=\partial_\beta\det P_0=0$ from
RPC6, without first assuming an inverse exists. At the origin the
literal Gamma ray integrals give the diagonal matrix $G_u$ defined in
RPC8, whose determinant is nonzero. Hence $P_0$ is invertible at every
point of $\mathbb C^2$, justifying the inverses throughout this family.

The matrix identity $M^{\mathsf T}J+JM=0$ also gives
$(M^3)^{\mathsf T}J+JM^3=0$. Direct multiplication of RPC5--6 gives
\[
 \partial_\beta J+\frac13(N^{\mathsf T}J+JN)=0.
 \tag{RPC7}
\]
In particular $\partial_\beta J$ has entries $(2,3)=1$, $(3,2)=-1$
when indices begin at zero, while $N^{\mathsf T}J+JN$ has the
opposite threefold entries. Differentiate the inverse matrices in
$P_0^{-\mathsf T}JP_0^{-1}$. RPC6--7 show that both its $\beta$ and
its $\eta$ derivatives vanish. This matrix is consequently constant
on the connected coefficient space $\mathbb C^2$.

Define the unchanged positive-Gamma ray constants
\[
 g_a=5^{a/5-1}e^{\pi ia/5}\Gamma(a/5),\qquad
 G=\operatorname{diag}(g_1,g_2,g_3,g_4),\quad
 G_u=\operatorname{diag}(u^{a/5}g_a)_{a=1}^4.
 \tag{RPC8}
\]
At $\beta=\eta=0$, literal substitution in each original ray integral
gives $P_0=G_u$. Every $g_a$ is nonzero, since its Gamma factor is
the integral of a strictly positive function. We have therefore proved
the complete flat-pairing identity
\[
 \begin{split}
 \Omega_0&=P_0^{-\mathsf T}JP_0^{-1}
       =G_u^{-\mathsf T}J\big|_{\beta=0}G_u^{-1},\\
 (\Omega_0)_{14}&=-\frac1{u g_1g_4},\qquad
 (\Omega_0)_{23}=\frac1{u g_2g_3},\qquad
 (\Omega_0)_{ba}=-(\Omega_0)_{ab},\\
 \Omega&=P^{-\mathsf T}JP^{-1}=e^{-2A/u}\Omega_0.
 \end{split}
 \tag{RPC9}
\]
All entries not listed are zero; the indices in RPC9 begin at one.
This formulation retains the Gamma constants directly and requires
no unevaluated normalization of a period or source mass. For
$D=\operatorname{diag}(\zeta,\zeta^2,\zeta^3,\zeta^4)$, the two
paired sums of weights are five. It follows entry by entry that
\[
 D^{\mathsf T}\Omega D=\Omega,
 \qquad D^{\mathsf T}\Omega_0D=\Omega_0.
 \tag{RPC10}
\]

\subsection{The full arithmetic unit and the exact quadratic form}

Retain the original $v_h=2\xi/h$ and its complete simple-root unit
values. Reflection and conjugation of this original function give
\[
 v_h(\alpha_1)=v_h(\alpha_4)=a,\qquad
 v_h(\alpha_2)=v_h(\alpha_3)=\overline a,\qquad a\ne0.
 \tag{RPC11}
\]
Indeed both numerator and denominator are invariant under $s\mapsto1-s$
and conjugate under $s\mapsto\overline s$, first away from the roots
and then by the removable values. The complete zero orders give
nonzero values at these simple roots. Write
\[
 c_3=\frac{a^{-2}-\overline a^{-2}}{4i\delta\gamma},\qquad
 c_1=\frac{a^{-2}+\overline a^{-2}}2-(\delta^2-\gamma^2)c_3,
 \quad \mathfrak r(w)=c_1w+c_3w^3.
 \tag{RPC12}
\]
Both coefficients are real, and $(c_1,c_3)\ne(0,0)$. Directly,
$c_1+c_3r_1^2=a^{-2}$ and $c_1+c_3r_2^2=\overline a^{-2}$;
the same equalities hold at the respective opposite roots. Thus
$\mathfrak r(M)=M\,v_h(M+1/2)^{-2}$ in the original finite algebra, with the
entire function evaluated by its complete original arithmetic unit.
RPC12 is an exact interpolation of this specific multiplier.

For $Q_0(t)=t_1t_4-t_2t_3$, let $J_{\rm prim}$ be the original
Lagrange coordinate isomorphism and set
\[
 B=F^{-1}\Pi UJ_{\rm prim},\qquad Q(x)=Q_0(B^{-1}x),\qquad
 \mathcal H=P\mathfrak r(M)P^{-1}.
 \tag{RPC13}
\]
If $p=P^{-1}x$, the primary components of $B^{-1}x$ are exactly
$p(r_1)/a,p(r_2)/\overline a,p(r_3)/\overline a,p(r_4)/a$.
Since $H'(r)=4r(r^2+\gamma^2-\delta^2)$, its values give
$r/H'(r)=1/(8i\delta\gamma)$ on roots one and four, and
$r/H'(r)=-1/(8i\delta\gamma)$ on roots two and three. Inserting
these four values into RPC4 proves the exact identity
\[
 \boxed{\quad Q(x)=4i\delta\gamma\,
       p^{\mathsf T}\mathfrak r(M)^{\mathsf T}Jp
       =4i\delta\gamma\,x^{\mathsf T}\mathcal H^{\mathsf T}\Omega x.\quad}
 \tag{RPC14}
\]
There is no conjugate transpose in this quadratic identity. Because
$\mathfrak r$ is odd, $\mathfrak r(M)^{\mathsf T}J+J\mathfrak r(M)=0$; hence the matrices in
RPC14 are symmetric. This also proves
$\mathcal H^{\mathsf T}\Omega+\Omega\mathcal H=0$.
The four eigenvalues of $\mathfrak r(M)$ are exactly
$r_1/a^2,\overline{r_1/a^2},-\overline{r_1/a^2},-r_1/a^2$.
They are nonzero, although they need not be distinct. Both $B$ and
$Q_0$ are invertible as linear and quadratic matrices, respectively,
so $Q$ always has rank four, including the possible eigenvalue collisions.

By RPC10, substituting $Dx$ into RPC14 transports $\mathcal H$ to
$D^{-1}\mathcal H D$. The resulting symmetric quadratic matrices
are equal if and only if the corresponding $\mathcal H$ matrices
are equal, because $\Omega$ is invertible. Consequently
\[
 \boxed{\quad Q(Dx)=Q(x)\quad\Longleftrightarrow\quad
 [\mathcal H,D]=0\quad\Longleftrightarrow\quad
 \mathcal H_{ab}=0\text{ for every }a\ne b.\quad}
 \tag{RPC15}
\]
This is an exact test in the actual original period matrix and unit.
The final equivalence uses the four distinct diagonal entries of $D$.
It makes no assumption that their off-diagonal tests are independent.

For completeness, a rank-four quadratic semi-invariant of cyclic
weight $j\ne0$ is impossible: its coefficient in row $j$ could pair
only with the missing weight zero, so that row is zero. Thus any
semi-invariant rank-four quadratic must have weight zero. Since five
is prime, the projective orbit of $Q$ under $D$ has size one or five,
and size one is exactly RPC15. Every failure of any displayed
off-diagonal test proves size five.

\subsection{An explicit original-parameter radius forcing orbit five}

Here we prove that RPC15 fails for all sufficiently large $|u|$,
with a completely specified finite radius for the given original
quartet and unit. The argument evaluates the original periods; it
does not choose or insert a replacement amplitude.

Set $b=\beta u^{-2/5}$, $c=\eta u^{-4/5}$, and let $T(b,c)$
be the Fourier period matrix for
$z^5/5+bz^3/3+cz$ at $u=1$, in columns $1,z,z^2,z^3$ and on
the five rays of angles $(\pi+2\pi j)/5$. Substitution
$w=u^{1/5}z$, with exactly the original branch, gives
\[
 P_0=T(b,c)S_u,\quad S_u=\operatorname{diag}(u^{a/5})_{a=1}^4,
 \quad S_uMS_u^{-1}=u^{1/5}M(b,c),\quad T(0,0)=G.
 \tag{RPC16}
\]
Every column includes its original differential factor. The matrix
$M(b,c)$ is RPC5 with $\beta,\eta$ replaced by $b,c$; in particular
$N_0:=M(0,0)$ has its three subdiagonal entries equal to one.

The following positive constants are explicit finite real integrals
or finite expressions in the original Gamma constants:
\[
 \begin{split}
 E_j&=\frac85\int_0^\infty r^j(r^3/3+r)
                    e^{-r^5/5+r^3/3+r}\,dr\quad(0\le j\le3),\\
 C_T&=2\left(\sum_{j=0}^3E_j^2\right)^{1/2},\qquad
 g_+=\max_a|g_a|,\quad g_-^{-1}=\max_a|g_a|^{-1},\\
 \kappa&=\max_{1\le a<4}|g_{a+1}/g_a|,\qquad
 K_0=GN_0G^{-1},\\
 C_K&=4C_Tg_-^{-1}+2g_+g_-^{-1}
                       +2g_+g_-^{-2}C_T,\\
 C_{K,3}&=C_K(3\kappa^2+3\kappa C_K+C_K^2),\qquad
 \epsilon_0=\min\{1,(2C_Tg_-^{-1})^{-1}\}.
 \end{split}
 \tag{RPC17}
\]
Here $g_-^{-2}$ means $(\max_a|g_a|^{-1})^2$. The negative fifth
degree exponent proves convergence of every integral. All norms in
the following estimates are the Euclidean operator norm, with the
Frobenius norm used explicitly as an upper bound where stated.

Write $\epsilon=|b|+|c|$. When $|b|,|c|\le1$, on each ray
\[
 |e^{bz^3/3+cz}-1|
 \le \epsilon(r^3/3+r)e^{r^3/3+r}.
\]
Each contour has two rays, and each entry of $F^{-1}$ has absolute
value $1/5$ with four summands. The entry in column $j$ of
$T-G$ is therefore at most $E_j\epsilon$. Summing the squares
over all four rows proves
\[
 \|T-G\|\le\|T-G\|_{\rm F}\le C_T\epsilon.
 \tag{RPC18}
\]
For $\epsilon\le\epsilon_0$, the convergent inverse series for
$I+G^{-1}(T-G)$ gives
$\|T^{-1}\|\le2g_-^{-1}$ and
$\|T^{-1}-G^{-1}\|\le2g_-^{-2}C_T\epsilon$.
Also $\|M(b,c)-N_0\|\le\epsilon$, $\|M(b,c)\|\le2$,
and $\|N_0\|=1$. Subtract the products
$K:=T M(b,c)T^{-1}$ and $K_0$ in three terms, changing the left
factor, then the middle factor, then the inverse factor. The preceding
inequalities give, with the complete RPC17 constants,
\[
 \|K-K_0\|\le C_K\epsilon,\quad
 \|K\|\le\kappa+C_K\epsilon,\quad
 \|K^3-K_0^3\|\le C_{K,3}\epsilon.
 \tag{RPC19}
\]
For the last assertion use the exact noncommutative identity
$K^3-K_0^3=(K-K_0)K^2+K_0(K-K_0)K+K_0^2(K-K_0)$.
Its norm is at most $C_K\epsilon$ times
$(\kappa+C_K)^2+\kappa(\kappa+C_K)+\kappa^2$, as claimed.

The original multiplication and full unit in RPC13 now give exactly
\[
 \mathcal H=c_1u^{1/5}K+c_3u^{3/5}K^3,
 \qquad (K_0^3)_{41}=g_4/g_1,\quad (K_0)_{21}=g_2/g_1.
 \tag{RPC20}
\]
The scalar $e^{A/u}$ cancels from this conjugation, while it remains
in the quadratic form through RPC9. Put $B_*=\beta+\eta>0$.
For $|u|\ge1$, $\epsilon\le B_*|u|^{-2/5}$.

If $c_3\ne0$, define the completely specified radius
\[
 \begin{split}
 T_*&=\max\left\{1,\frac{B_*}{\epsilon_0},
 \frac{2\bigl(|c_3|C_{K,3}B_*+|c_1|(\kappa+C_K)\bigr)}
 {|c_3g_4/g_1|}\right\},\qquad R_*=T_*^{5/2}.
 \end{split}
 \tag{RPC21}
\]
For every $u$ on its original branch with $|u|\ge R_*$,
RPC19--20 imply
\[
 \left|u^{-3/5}\mathcal H_{41}-c_3g_4/g_1\right|
 \le |u|^{-2/5}\bigl(|c_3|C_{K,3}B_*
                           +|c_1|(\kappa+C_K)\bigr)
 \le\tfrac12|c_3g_4/g_1|.
 \tag{RPC22}
\]
In particular $\mathcal H_{41}\ne0$. In the quadratic matrix
RPC14 this is the nonzero $(1,1)$ entry
$4i\delta\gamma\,\mathcal H_{41}\Omega_{41}$, whose cyclic
weight is two.

If $c_3=0$, then $c_1\ne0$. In this case define
\[
 T_* =\max\left\{1,\frac{B_*}{\epsilon_0},
                   \frac{2C_KB_*}{|g_2/g_1|}\right\},
 \qquad R_*=T_*^{5/2}.
 \tag{RPC23}
\]
For every $|u|\ge R_*$, the same estimates give
\[
 \left|u^{-1/5}\mathcal H_{21}-c_1g_2/g_1\right|
       \le |c_1|C_KB_*|u|^{-2/5}
       \le\tfrac12|c_1g_2/g_1|.
 \tag{RPC24}
\]
Thus $\mathcal H_{21}\ne0$. The $(1,3)$ entry of the quadratic
matrix is then $4i\delta\gamma\mathcal H_{21}\Omega_{23}$,
which has cyclic weight four; symmetry retains its $(3,1)$ partner.
Both cases prove the unconditional conclusion for the actual original
parameters and the explicit radius belonging to those parameters:
\[
 \boxed{\quad |u|\ge R_*\quad\Longrightarrow\quad
       \#\{[Q(D^jx)]:0\le j<5\}=5.\quad}
 \tag{RPC25}
\]
The bracket denotes the projective class of the quadratic polynomial.
RPC21 or RPC23 supplies its radius for every nonzero original unit;
neither case assumes a condition on an uncomputed period coefficient.

Finally, the original periods are holomorphic on the connected
logarithm cover of $\mathbb C^*$, and on each branch domain: locally the contours may
be held in common decay sectors and differentiated under the dominated
integral, and contour deformation identifies these local definitions.
Their determinant never vanishes. Hence all entries of $\mathcal H$
are holomorphic there. RPC22 or RPC24 proves that its respective
decisive entry is not identically zero on the logarithm cover and,
by the identity theorem on that connected cover, on any open branch domain. The zeros
of that one entry are isolated, by its convergent local power series.
The exact exceptional set
\[
 \mathcal E=\{u:\mathcal H_{ab}(u)=0\text{ for all }a\ne b\}
 \tag{RPC26}
\]
is a subset of those isolated zeros and satisfies $|u|<R_*$.
Thus it has no accumulation inside the branch domain; accumulation
at $u=0$ or a branch boundary has not been excluded. RPC15 is the
full calculation specifying this set. No assertion that it is empty
at arbitrary finite complex $u$ is made or used.

\end{document}
